\section{结论}
\label{sec:conclusion}

本文面向强约束三维低空无人机路径规划问题，提出了约束可行性场驱动的自适应进化搜索方法 CVF-AE。该方法通过约束感知初始化、局部 CVF 候选引导、状态自适应触发和稀疏可行性保持，将约束违反信息用于低开销的方向修正，而不是在每一代对全种群进行强制更新。

三类合成低空场景实验表明，CVF-AE 在所测试场景中保持稳定可行性，并在主对比中取得最佳平均排名。统计检验和轨迹复核进一步说明，本文方法的优势主要体现在可行率、轨迹合理性和评价开销之间的均衡，而不是保证每个场景都获得最低标量适应度。消融结果验证了 CVF、约束感知初始化和稀疏保持机制在强约束环境中的互补作用；评价开销结果则表明，CVF 分支以稀疏方式触发，额外评价次数保持在约 $1\%$ 量级。

总体而言，CVF-AE 为强约束低空路径规划提供了一种可解释、可控开销的工程型进化优化框架。后续研究将围绕真实地形环境、动态障碍约束以及规划--控制闭环验证进一步扩展。
